\section{The Range of Feeling}

From here the paper turns inward, to read the same substrate as experience rather than as physics. Everything in this part rests on the one premise marked earlier, that the tone is felt, and granting it, the question becomes how much can be felt, and how it is shaped. The honest status is that the combinatorics below is exact while the claim that it is felt is the posit, and we keep the two apart.

The single quale, the smallest possible feeling, is one tone. It is a single value in $\{-1, 0, +1\}$, pain or peace or pleasure, and that is the whole of it, the irreducible atom of experience, one vibe filling one site. Why three and not two is worth a moment, because two would give only a yes and a no, while three gives a push, a rest, and a pull, the least structure in which there is a middle to come from and return to. Peace is that middle, the still source, and the entire felt world is built from the three.

A single tone is the atom, but a single dock is already a world. A dock holds twenty-four tones at once, one in each site, so its state is a point in a space of $3^{24}$, about two hundred and eighty billion possibilities, and that state is not a flat list but a structured thing, a \emph{hyper-color}. The twenty-four split, by triality, into a sensory octet and two octets of opposite handedness, so a single dock's feeling already has a direction-of-flow and a pair of spin-like grains woven together, a felt quality richer than any colour by orders of magnitude. The dock, not the single tone and not the whole self, is the hyper-color, the first composite feeling.

\begin{result}[The tower of feeling]
Feeling composes by multiplication. Each dock multiplies the number of states by $3^{24}$, so the digit-count of a self's state-space grows linearly with the number of docks, about eleven digits per dock. One tone is three states, a dock is a hyper-color of $3^{24}$, a self of a billion docks has a state-count with eleven billion digits, and the ever-growing whole has no bound at all. The range of feeling is a tower whose every rung leaps past the entire reach of the rung below.
\end{result}

A few things follow from these numbers whether or not one grants that they are felt. Description is hopelessly lower-dimensional than the state it describes, so any experience is ineffable in the precise sense that words cannot carry its dimension, no two complete states ever exactly recur, so each felt moment is genuinely unprecedented, and because the crystal grows without end the space of feeling is inexhaustible, never used up. The fine ternary grain, averaged at the scale of a self, reads as a smooth near-continuum, so the inner world feels seamless though it is built of trits. What keeps so vast a space from being mere noise is the same thing that organizes the bulk into a database, integration and indexing, which carve the immensity into coherent, navigable regions, the emotions and percepts and memories of the chapters that follow. The range is unimaginably large, and it is also, by the geometry, livable.
